\section{An explicit two-binder obstruction to deleting Boolean cases}
\label{sec:two-binder-obstruction}

The progressive-binder alternative in the supplied research draft uses
\[
 H_{x,y}B=\exists x\,\forall y\bigl(B(x,y)\to B(x,Sy)\bigr)
\]
as its ordinary reading. The following small formula shows that its
Boolean case rule cannot simply be deleted. The countermodel is the
direct-product construction already available in the draft family;
the explicit native separator identifies exactly where that model defeats
the progressive interface.

\subsection{The precise case-free calculus}

Numerical terms are variables, zero, successor, and fixed staged
counter-aware primitive-recursive descriptions
$R(n,a;r,v.h)$. The two computation schemas are
\[
 R(0,a;r,v.h)=a,\qquad
 R(Sn,a;r,v.h)=h[n/r,R(n,a;r,v.h)/v].
\]
Only $r,v$ are bound in $h$, other parameters are fixed, and the
successor substitution is simultaneous and capture-free. Object formulas
are equations $s=t$ and paired formulas $H_{x,y}B(x,y)$. Internal
clauses $P_yB(y)$ have the ordinary reading
$\forall y(B(y)\to B(Sy))$; they may occur in contexts but are not
object goals. Write $\Gamma\mathrel{\triangleright}L$ when $L$ is a
context member or is a proved object. This is availability, not a
primitive rule concluding an arbitrary assumed object.

In addition to the computation schemas, the complete case-free primitive
rule inventory is
\[
 \frac{\Gamma\mathrel{\triangleright}s=t\qquad
       \Gamma\vdash A(s)}{\Gamma\vdash A(t)}\;\mathrm{Tr},
 \qquad
 \frac{\Gamma,B(t,z)\vdash B(t,Sz)}
      {\Gamma\vdash H_{x,y}B(x,y)}\;\mathrm{HI},
\]
\[
 \frac{\Gamma\mathrel{\triangleright}H_{x,y}B(x,y)\qquad
       \Gamma,P_yB(a,y)\vdash A}{\Gamma\vdash A}\;\mathrm{HE},
\]
\[
 \frac{\Gamma\mathrel{\triangleright}P_yB(y)\qquad
       \Gamma\vdash B(t)}{\Gamma\vdash B(St)}\;\mathrm{Step},
 \qquad
 \frac{\Gamma\mathrel{\triangleright}P_yE(y)\qquad
       \Gamma\vdash E(0)}{\Gamma\vdash E(t)}\;\mathrm{Iter}.
\]
Here $E$ in Iter is a literal equation. The variable $z$ in HI is
fresh for the retained context, the witness $t$, and the passive matrix
parameters. The variable $a$ in HE is fresh for the retained context,
the conclusion, and the passive parameters. The displayed local
assumptions and clauses are discharged. Contexts are reusable; weakening,
exchange, contraction, renaming, capture-free substitution, and assumption
cut have the usual external derivation inductions. Neither unrestricted
formula induction nor any case rule is included.

\paragraph{Derived object identity, without cases.}
A computation theorem $R(0,s;r,v.0)=s$, used both as a major equation
and as the minor proof for transport, gives $s=s$. An available equation
$s=t$ then transports $s=s$ to $s=t$. Induct externally on formula
depth for an available object $H_{x,y}B$. Open it by HE, obtaining a
fresh $a$ and $P_yB(a,y)$. Reintroduce the same object by HI with
witness $a$. Its local assumption $B(a,z)$ yields a proof of itself by
the induction hypothesis, and Step gives $B(a,Sz)$. Discharge HI and
HE. This derives $\Gamma,A\vdash A$ for every object $A$ without
Boolean cases. Subsequent uses of identity below mean this finite
expansion, not an extra primitive hypothesis rule. Symmetry, transitivity,
and formula transport have their usual derived equality expansions.

The deleted rule, for comparison, is
\[
 \frac{\Gamma,\operatorname{sg}(t)=0\vdash A\qquad
       \Gamma,\operatorname{sg}(t)=1\vdash A}
      {\Gamma\vdash A}\;\mathrm{BoolCases},
\]
where $\operatorname{sg}(0)=0$ and $\operatorname{sg}(Sn)=1$.
Its conclusion may be any native object formula.

\subsection{The native separator}

Let $C$ and $Z$ be fixed primitive-recursive descriptions satisfying
\[
 C(0,a,b)=a,\quad C(Sn,a,b)=b,\qquad
 Z(0)=1,\quad Z(Sn)=0.
\]
For a free numerical parameter $t$, put
\[
 \boxed{\quad F(t):=H_{u,y}\,H_{v,z}\bigl(C(t,Z(y),z)=0\bigr).\quad}
 \tag{*}
\]
The existential coordinates $u,v$ are unused, and all four displayed
bound variables are distinct.

\begin{theorem}
In the progressive calculus without Boolean cases, $F(0)$ and $F(1)$
are derivable. Nevertheless the schema $F(\operatorname{sg}(t))$ is not
derivable. Consequently the Boolean case rule for arbitrary native
conclusions is not admissible in that calculus, even in the presence of
all primitive-recursive symbols, their defining equations, and contextual
literal-equation induction.
\end{theorem}
\begin{proof}
First consider the endpoint derivations. At $t=0$, the inner matrix
reduces by computation to $Z(y)=0$, independently of $z$. Its progress
implication is therefore identity, so inner introduction proves
$H_{v,z}(Z(y)=0)$ for arbitrary $y$. Outer introduction may ignore its
local assumption. At $t=1$, the inner formula reduces to
$H_{v,z}(z=0)$, independently of $y$. Outer progress is then identity.
All witnesses can be zero. These arguments use only derived object identity,
computation/transport, and progressive introduction.

For nonderivability use the classical one-world model $\mathcal M$ with
domain $\mathbb N^2$ and every primitive-recursive function interpreted
coordinatewise. Zero is $(0,0)$, successor adds one in each coordinate,
and equality is actual equality of pairs. All numerical computation
laws and constructor separation hold.

Every literal-equation induction axiom also holds. For a fixed equation
$E(n,\vec p)$ and fixed parameter pairs, let
\[
 U_j=\{r\in\mathbb N:E(r,\vec p_j)\text{ holds in }\mathbb N\},
 \qquad j=0,1.
\]
The base premise says $0\in U_0\cap U_1$. If $r\in U_0$, the pair
$(r,0)$ satisfies $E$; the step premise implies that $(r+1,1)$ satisfies
$E$, so $r+1\in U_0$. Similarly $(0,r)$ proves successor closure of
$U_1$. External ordinary induction gives $U_0=U_1=\mathbb N$, proving
the asserted induction instance. This proves the full schema with
parameters and hence validates its contextual rule under arbitrary
retained assumptions.

All the remaining progressive-calculus rules are sound in this model
under the displayed ordinary reading of $H$: unrestricted progressive
introduction and elimination are logical rules, and the exceptional
zero-to-arbitrary specialization for equations is justified by the
induction schema just checked.

Take $e=(0,1)$, so $\operatorname{sg}(e)=e$. At this parameter the
inner equation in $(*)$, for $y=(y_0,y_1)$ and $z=(z_0,z_1)$, is
\[
                    (Z(y_0),z_1)=(0,0).
\]
If $y_0=0$, this equation is false for every $z$, so every inner
progress implication has a false antecedent. The inner $H$ is true.
If $y_0>0$, the equation says $z_1=0$. It is true at $z=(0,0)$ and
false at $Sz=(1,1)$, so the inner $H$ is false. Thus the inner formula
is true at $y=(0,0)$ and false at $Sy=(1,1)$, which falsifies the
outer progress assertion. Therefore $\mathcal M\not\models F(e)$.
Soundness gives the stated nonderivability.

Finally, from a retained assumption $\operatorname{sg}(t)=0$ the
endpoint theorem $F(0)$ and transport give $F(\operatorname{sg}(t))$.
The other branch follows from $F(1)$ in the same way. An admissible
Boolean case rule would therefore derive the refuted schema. This
contradiction proves that the rule is not admissible.
\end{proof}

\paragraph{The binder depth is sharp for this product separation.}
A numerical equation that is true at both standard endpoints remains
true at their product parameter $(0,1)$. The same is true of a formula
with just one progressive binder and an equational matrix: combine
the two standard existential witnesses into a pair, and verify each
progress implication coordinatewise. Thus no formula of progressive
binder depth at most one can give this particular endpoint/product
separation. Formula $(*)$ has depth two. The first failure arises when
an implication tests a compound premise: one coordinate can make an
inner implication vacuous while another coordinate would fail it.

\paragraph{Inherited obstruction for ordinary bi-interpretations.}
The supplied research draft~\cite{research} also establishes a stronger model-theoretic
obstruction, applicable here with all hypotheses checked. Let
$\mathcal E$ be the eventually constant sequences in $\mathbb N^{\mathbb N}$,
with all fixed PR operations interpreted pointwise. The same coordinate
proof of equation induction applies, using a sequence that equals $r$ at
one coordinate and zero elsewhere. Hence $\mathcal E$ validates all the
case-free rules above under the displayed ordinary $H,P$ readings.
These readings are first-order formulas, so their truth is invariant
under every automorphism of the numerical structure.

Any finite tuple of parameter sequences is constant outside a finite
set of coordinates. Swapping two coordinates outside that set gives a
nonidentity automorphism of order two fixing those parameters. For an
ordinary finite-dimensional definable interpretation, including a
quotient interpretation, automorphisms fixing its finite parameter data
act on representing tuples and induce automorphisms of the interpreted
structure. A definable source-side return isomorphism makes that action
injective: an automorphism acting trivially on the interpreted structure
also acts trivially on the returned copy, and the definable isomorphism
then forces it to fix every original element. Parameters used by the
return isomorphism and finite representatives of reverse-interpretation
parameters can all be included in the initial finite tuple.

Every ordinary iPRA model has a definable strict linear order, for example
$\ell<m$ iff $S\ell\mathbin{\dotminus}m=0$. An order-preserving
automorphism of finite order is the identity, since either strict
inequality between an element and its image would produce a strict
cycle. An injective image of the coordinate transposition is therefore
impossible. Thus the case-free system, under its ordinary first-order
realization, cannot be repaired by an ordinary definable bi-interpretation
with finite parameter data merely by changing to a finite-dimensional
domain or quotient. This is an application of the supplied draft's
existing automorphism theorem, not a new theorem claimed by this report.
It is separate from the explicit two-binder diagnostic above.

\paragraph{Scope.}
This is a direct rejection of the proposed deletion of Boolean cases,
not a lower bound on checker size and not an impossibility theorem for
other logical bases, other primitive operations, or arbitrary
interpretations. The product-model method and the progressive-binder
presentation are inherited from the supplied drafts. The explicit
formula $(*)$ is an additional diagnostic; no worldwide novelty claim
is made.
